% Demonstrative expansion DF003--DF007. The integral source is unchanged.
% Common owner of the Lean files cited in these comments:
% output/PAQUETE_ARTICULO_I_INTEGRACION_ESTADO_CAMPOS_20260919/.
% Causal placement of the first section: after the coinductive regional
% readers and before their registers are used as inputs to K selection.
% Causal placement of the second section: after the marked
% Paley--Witt--Golay--Leech construction and before reconstruction of the VOA.
% Both expository cuts preserve the same APP--TRIT--TPK origin.

\section[Regional generation and refinement]{Integral compatibility of regional generation, cylindrical refinement, and incidence}
\label{act21:sec:df-lectores-cilindros}

Regional prolongation produces ordered sequences of blocks together with
their reading depths. The Archimedean projection and ternary incidence
are calculated on these same sequences. The connection between the two
realizations can be formulated entirely through integer operations:
selection by rational intervals, Euclidean division, lifting of six
trits, carry reading, and the Witt transformation. This formulation makes
it possible to trace each coefficient from its generating block to its
subsequent use.

The construction uses the three regional readers already obtained from the
APP--TRIT--TPK state. We denote them by the indices
\(c\in\{\mathrm P,\mathrm E,\mathrm F\}\), corresponding, respectively,
to the closure, propagation, and autoscaling regions. Their real realization
will be written as \(v_c\). The subsequent identifications
\(v_{\mathrm P}=\pi_{\mathrm{HMT}}\),
\(v_{\mathrm E}=e_{\mathrm{HMT}}\), and
\(v_{\mathrm F}=\varphi_{\mathrm{HMT}}\) preserve this causal order.
The operation that emits a block inspects the rational intervals of the
regional reader; the real value enters the proof of correctness of that
operation.

\subsection{Positional emission through rational separation of cells}
\label{act21:subsec:df-emision-racional}

% Source: lean/biblioteca/RegionalPublicationComposition.lean,
% produced_brackets, eventual_publication, joint_eventual_publication,
% search_terminates, stopped_cell_correct, publications_compatible.

For each region, two rational sequences
\(\ell_c(j),u_c(j)\), constructed by its reader, are available and satisfy
\[
 \ell_c(j)\leq v_c\leq u_c(j).
\]
The cell-separation property proved for these readers states that, for
each integer \(s>0\), there exists \(J\) such that, for every \(j\geq J\),
\begin{equation}
 \lfloor s\ell_c(j)\rfloor
 =\lceil su_c(j)\rceil-1
 =\lfloor sv_c\rfloor.
 \label{act21:eq:df-separacion-celdas}
\end{equation}
Avoidance of rational boundaries by the regional values is the hypothesis
of the preceding theorems that guarantees this separation. Here we use
their operational conclusion together with the readers and intervals
that produce it.

\begin{definition}[Separation index and emitted prefix]
For \(s>0\), let
\[
 j_c(s)=\min\{j\in\mathbb N:
 \lfloor s\ell_c(j)\rfloor=\lceil su_c(j)\rceil-1\}.
\]
For an integer base \(b\geq2\) and a depth \(n\geq0\), define
\begin{equation}
 P_c(b,n)=\lfloor b^n\ell_c(j_c(b^n))\rfloor.
 \label{act21:eq:df-publicacion-prefijo}
\end{equation}
The prefix contains the integer part and the first \(n\) positions in base
\(b\); consequently, \(P_c(b,0)\) is the regional integer part.
\end{definition}

\begin{theorem}[Termination and correctness of emission]
\label{act21:thm:df-emision-correcta}
The index \(j_c(s)\) exists for every \(s>0\). The emitted prefixes satisfy
\begin{equation}
 P_c(b,n)=\lfloor b^n v_c\rfloor,
 \qquad
 P_c(b,n)\leq b^n v_c<P_c(b,n)+1.
 \label{act21:eq:df-prefijo-correcto}
\end{equation}
For each scale \(s\), a single sufficiency index \(J(s)\) can also be
chosen that is valid simultaneously for all three regions.
\end{theorem}

\begin{proof}
The eventual equality in \eqref{act21:eq:df-separacion-celdas} makes the set
defining \(j_c(s)\) nonempty. Its minimum exists by the well-ordering of
\(\mathbb N\). At any index where the equality holds, the rational lower
bound gives
\(\lfloor s\ell_c(j)\rfloor\leq\lfloor sv_c\rfloor\).
The upper bound, together with avoidance of a cell boundary, gives
\(\lfloor sv_c\rfloor\leq\lceil su_c(j)\rceil-1\).
When the two integer endpoints coincide, the intervening integer coincides
with them. This proves \eqref{act21:eq:df-prefijo-correcto}. For the joint
statement, it suffices to take the maximum of the three sufficiency indices
obtained at the same scale. The joint index controls all three emissions
simultaneously, while each region retains its own intervals.
\end{proof}

\begin{proposition}[Exact compatibility of prefixes]
\label{act21:prop:df-prefijos-compatibles}
For \(n,m\geq0\),
\begin{equation}
 \left\lfloor\frac{P_c(b,n+m)}{b^m}\right\rfloor=P_c(b,n).
 \label{act21:eq:df-truncamiento}
\end{equation}
In particular, the next block
\(d_c(b,n)=P_c(b,n+1)\bmod b\) satisfies
\begin{equation}
 P_c(b,n+1)=bP_c(b,n)+d_c(b,n),\qquad 0\leq d_c(b,n)<b.
 \label{act21:eq:df-paso-posicional}
\end{equation}
\end{proposition}

\begin{proof}
Write \(b^n v_c=N+\theta\), with
\(N=P_c(b,n)\) and \(0\leq\theta<1\). Then
\(P_c(b,n+m)=b^mN+\lfloor b^m\theta\rfloor\), and the second summand
belongs to \(\{0,\ldots,b^m-1\}\). Euclidean division by \(b^m\)
recovers \(N\). The specialization \(m=1\) yields
\eqref{act21:eq:df-paso-posicional} and its residual bound.
\end{proof}

\subsection{Regional blocks of six trits at arbitrary depth}
\label{act21:subsec:df-bloques-arbitrarios}

% Source: deltas/integracion/JointRegionalFrontier.lean,
% publish_base729_eq_base3, generated_block_from_longer_prefix,
% generated_block_eq_finite, generated_signature_eq_finite.

The integer equality \(729=3^6\) determines the grouping of six ternary
positions into a block. For \(t\geq0\), define
\begin{equation}
 u_c(t)=P_c(729,t+1)\bmod729,
 \qquad
 B_{c,j}(t)=
 \left\lfloor\frac{u_c(t)}{3^{5-j}}\right\rfloor\bmod3,
 \quad 0\leq j<6.
 \label{act21:eq:df-bloque-y-banda}
\end{equation}
Thus \(B_c(t)\) is an ordered word of six trits, and
\begin{equation}
 u_c(t)=\sum_{j=0}^{5}B_{c,j}(t)3^{5-j}.
 \label{act21:eq:df-elevacion-banda}
\end{equation}
The word and its integer lift are mutually recoverable representations
of the same emitted block.

\begin{theorem}[Stability of extraction at any depth]
\label{act21:thm:df-bloque-estable}
For every \(n\),
\[
 P_c(729,n)=P_c(3,6n).
\]
If \(t<n\), then
\begin{equation}
 u_c(t)=
 \left\lfloor\frac{P_c(729,n)}{729^{n-t-1}}\right\rfloor\bmod729.
 \label{act21:eq:df-bloque-prefijo-largo}
\end{equation}
Consequently, every extension of the prefix preserves all previously
emitted blocks and all signatures calculated from them.
\end{theorem}

\begin{proof}
The first equality follows because both prefixes are separated at the same
scale, \(729^n=3^{6n}\), and from
Theorem~\ref{act21:thm:df-emision-correcta}. For the second, apply
\eqref{act21:eq:df-truncamiento} at depths \(n\) and \(t+1\), and then take
the residue modulo \(729\). The ternary extraction in
\eqref{act21:eq:df-bloque-y-banda} is deterministic. Any explicit function of
those eighteen coordinates therefore retains the same result when the
prefix is extended.
\end{proof}

The construction is quantified over \(t\in\mathbb N\). A table
materializing one hundred blocks is a finite evaluation of this
construction; formula \eqref{act21:eq:df-bloque-prefijo-largo} determines its
compatibility with every subsequent extension exactly.

\begin{example}[First joint block]
Evaluation of the selected regional readers gives the words
\[
 B_{\mathrm P}(0)=010211,\qquad
 B_{\mathrm E}(0)=201101,\qquad
 B_{\mathrm F}(0)=121200.
\]
Their lifting by \eqref{act21:eq:df-elevacion-banda} produces
\[
 u_{\mathrm P}(0)=103,\qquad
 u_{\mathrm E}(0)=523,\qquad
 u_{\mathrm F}(0)=450.
\]
For example,
\(201101_3=2\cdot243+27+9+1=523\).
All six positions, including leading zeros, are retained in the band
representation. The leading zero in \(010211\) has a defined position
even though the decimal notation for \(103\) does not contain it.
\end{example}

\subsection{Transition registers and reconstruction of linear lifts}
\label{act21:subsec:df-registros-transicion}

% Sources: deltas/integracion/RegionalTransitionRecords.lean and
% GeneratedTransitionRecords.lean; lean/biblioteca/TPKLifts.lean.
% The reconstructed links are R(0)->R(1) and R(2)->R(3).

Over \(\mathbb F_3\), collect two consecutive emissions from each region
in the matrix
\begin{equation}
 R(t)=
 \begin{pmatrix}
 B_{\mathrm P}(t)\\ B_{\mathrm P}(t+1)\\
 B_{\mathrm E}(t)\\ B_{\mathrm E}(t+1)\\
 B_{\mathrm F}(t)\\ B_{\mathrm F}(t+1)
 \end{pmatrix}.
 \label{act21:eq:df-registro-matricial}
\end{equation}
This definition fixes both the regional order and the temporal order.
The first four registers are
\[
 X_0=R(0),\quad Y_0=R(1),\quad X_1=R(2),\quad Y_1=R(3).
\]
Their values, obtained by extraction from the regional prefixes, are
\begin{align*}
 X_0&=\begin{pmatrix}
0&1&0&2&1&1\\0&1&2&2&2&2\\2&0&1&1&0&1\\
1&2&1&2&2&1\\1&2&1&2&0&0\\1&1&2&2&0&2
\end{pmatrix},&
 Y_0&=\begin{pmatrix}
0&1&2&2&2&2\\0&1&0&2&1&1\\1&2&1&2&2&1\\
1&0&2&0&1&1\\1&1&2&2&0&2\\1&2&1&0&2&0
\end{pmatrix},\\[1ex]
 X_1&=\begin{pmatrix}
0&1&0&2&1&1\\0&0&2&1&1&1\\1&0&2&0&1&1\\
0&1&2&2&2&2\\1&2&1&0&2&0\\0&1&0&2&1&0
\end{pmatrix},&
 Y_1&=\begin{pmatrix}
0&0&2&1&1&1\\1&1&0&2&2&1\\0&1&2&2&2&2\\
1&0&2&0&1&1\\0&1&0&2&1&0\\0&1&0&2&0&0
\end{pmatrix}.
\end{align*}

\begin{theorem}[Unique reconstruction of the two initial links]
\label{act21:thm:df-elevaciones-reconstruidas}
The equations \(X_0L_0=Y_0\) and \(X_1L_1=Y_1\) each have a unique
solution in \(M_6(\mathbb F_3)\). These solutions are
\begin{align*}
 L_0&=\begin{pmatrix}
2&2&2&1&2&1\\2&1&2&2&1&1\\1&0&0&1&0&2\\
0&0&1&1&1&0\\2&2&2&0&2&1\\2&1&2&1&0&0
\end{pmatrix},&
 L_1&=\begin{pmatrix}
2&1&2&0&2&2\\1&2&1&1&2&0\\1&2&1&1&1&2\\
0&1&0&0&2&1\\2&0&2&1&0&1\\0&2&2&2&1&1
\end{pmatrix}.
\end{align*}
Both matrices are invertible.
\end{theorem}

\begin{proof}
The matrices
\begin{align*}
 A_0&=\begin{pmatrix}
1&1&2&0&1&2\\0&0&1&0&0&1\\2&1&2&1&0&1\\
0&2&0&1&0&2\\2&1&1&0&2&2\\2&1&1&1&1&2
\end{pmatrix},&
 A_1&=\begin{pmatrix}
1&0&1&2&0&0\\0&0&1&1&2&2\\0&1&2&0&1&1\\
1&1&0&2&0&0\\1&1&2&1&1&2\\1&0&0&0&0&2
\end{pmatrix}
\end{align*}
satisfy \(A_iX_i=X_iA_i=I\), by multiplication in \(\mathbb F_3\).
The solutions are therefore forced by \(L_i=A_iY_i\); multiplication
produces the two displayed matrices. Their inverses are
\begin{align*}
 L_0^{-1}&=\begin{pmatrix}
1&2&0&2&0&2\\2&0&2&2&0&0\\2&1&2&0&2&0\\
1&0&0&0&2&0\\0&2&1&1&2&0\\2&2&2&2&2&2
\end{pmatrix},&
 L_1^{-1}&=\begin{pmatrix}
2&0&2&1&2&1\\2&1&0&0&2&0\\2&0&2&2&0&2\\
2&0&0&1&1&0\\1&1&1&1&1&0\\2&0&1&2&2&0
\end{pmatrix}.
\end{align*}
The bilateral verification \(L_iL_i^{-1}=L_i^{-1}L_i=I\) proves exact
recovery of each transformed word. All products use residues \(0,1,2\)
with their interpretation in \(\mathbb F_3\).
\end{proof}

The causal reading of this result is precise: the regional emissions
produce \(R(t)\); the registers \(R(0),\ldots,R(3)\) determine the two
lifts above; their inverses recover the transformed words. Continuity
for every \(t\) belongs to the coinductive readers and
Theorem~\ref{act21:thm:df-bloque-estable}. The matrix reconstruction proved here
has the two links specified in its statement as its domain.

\subsection{Simultaneous refinement in ternary and decimal bases}
\label{act21:subsec:df-cilindros-mixtos}

% Source: deltas/integracion/RegionalCylinderTransition.lean.

A ternary emission of six trits and a decimal emission of three digits
are readings in different bases. To preserve their relationship, record
the state
\[
 s=(n,k,N,D),
\]
where \(N\) is a prefix at depth \(n\) in base \(729\), and \(D\) is
a prefix at depth \(k\) in base \(1000\). Their Archimedean cylinders are
\begin{equation}
 I_{729}(s)=\left[\frac{N}{729^n},\frac{N+1}{729^n}\right),\qquad
 I_{1000}(s)=\left[\frac{D}{1000^k},\frac{D+1}{1000^k}\right).
 \label{act21:eq:df-cilindros}
\end{equation}
The half-open upper endpoint assigns each value to a unique cell at a
given base and depth.

\begin{definition}[Integer compatibility defects]
Define
\begin{align}
 A(s)&=N1000^k-D729^n,\label{act21:eq:df-defecto-inferior}\\
 B(s)&=(N+1)1000^k-D729^n=A(s)+1000^k.
 \label{act21:eq:df-defecto-superior}
\end{align}
The state is compatible when
\begin{equation}
 A(s)<729^n,\qquad B(s)>0.
 \label{act21:eq:df-compatibilidad}
\end{equation}
\end{definition}

\begin{lemma}[Integer criterion for intersection]
The cylinders in \eqref{act21:eq:df-cilindros} have a nonempty intersection if
and only if \eqref{act21:eq:df-compatibilidad} holds.
\end{lemma}

\begin{proof}
Two half-open intervals \([a,b)\) and \([c,d)\), with \(a<b\) and
\(c<d\), intersect exactly when \(a<d\) and \(c<b\). Applying this
condition to \eqref{act21:eq:df-cilindros} and multiplying by the positive
denominators gives
\(N1000^k<(D+1)729^n\) and
\(D729^n<(N+1)1000^k\). These are the two inequalities in
\eqref{act21:eq:df-compatibilidad}.
\end{proof}

\begin{definition}[Update of depths and prefixes]
For \(0\leq u<729\) and \(0\leq d<1000\), define
\begin{align*}
 T_u(n,k,N,D)&=(n+1,k,729N+u,D),\\
 T_{u,d}(n,k,N,D)&=(n+1,k+1,729N+u,1000D+d).
\end{align*}
The first operation refines only the ternary cylinder. The second refines
both readings and preserves both depths in the state.
\end{definition}

\begin{proposition}[Exact boundary recurrences]
\label{act21:prop:df-rec-frontera}
The updates above satisfy
\begin{align}
 A(T_u s)&=729A(s)+u1000^k,\label{act21:eq:df-rec-ternaria}\\
 A(T_{u,d}s)&=729000A(s)+u1000^{k+1}-d729^{n+1}.
 \label{act21:eq:df-rec-conjunta}
\end{align}
\end{proposition}

\begin{proof}
For the first identity,
\[
 (729N+u)1000^k-D729^{n+1}
 =729(N1000^k-D729^n)+u1000^k.
\]
For the second,
\begin{align*}
 &(729N+u)1000^{k+1}-(1000D+d)729^{n+1}\\
 &\qquad=729000(N1000^k-D729^n)
      +u1000^{k+1}-d729^{n+1}.
\end{align*}
Both are integer identities preceding any real approximation.
\end{proof}

\begin{theorem}[Refinement of the generated regional cylinders]
\label{act21:thm:df-refinamiento-generado}
Let
\[
 s_c(n,k)=(n,k,P_c(729,n),P_c(1000,k)).
\]
Then \(s_c(n,k)\) is compatible for all \(n,k\). Moreover,
\begin{align*}
 T_{u_c(n)}s_c(n,k)&=s_c(n+1,k),\\
 T_{u_c(n),d_c(1000,k)}s_c(n,k)&=s_c(n+1,k+1).
\end{align*}
Euclidean division of the refined prefixes recovers the preceding
prefixes exactly. For every \(a,b\geq0\),
\[
 \frac{P_c(729,n+a)-P_c(729,n+a)\bmod729^a}{729^a}
 =P_c(729,n),
\]
and the analogous identity holds in base \(1000\).
\end{theorem}

\begin{proof}
Theorem~\ref{act21:thm:df-emision-correcta} places \(v_c\) simultaneously in
both intervals in \eqref{act21:eq:df-cilindros}; the intersection criterion
gives compatibility. The update identities are
\eqref{act21:eq:df-paso-posicional} in each base. Recovery follows from
\eqref{act21:eq:df-truncamiento}. Thus the defect recurrence and prefix
recovery are consequences of the same update.
\end{proof}

\begin{example}[First pair of regional cylinders]
For \(n=k=1\), evaluation of the three regions gives
\[
\begin{array}{c|rr|rr}
 c&N&D&A&B\\\hline
 \mathrm P&2290&3141&211&1211\\
 \mathrm E&1981&2718&-422&578\\
 \mathrm F&1179&1618&-522&478
\end{array}
\]
In all three rows, \(A<729\) and \(B>0\). The first prefix column
decomposes as \(2290=3\cdot729+103\),
\(1981=2\cdot729+523\), and \(1179=729+450\): exactly the previously
emitted six-trit blocks reappear. The different signs of \(A\) express
the relative positions of the two lower endpoints without altering the
regional value's membership in both intervals.
\end{example}

\subsection{Selection of the next block by the regional reader}

Compatibility of two intervals expresses a relationship between readings.
Determining the next block additionally uses the rational interval
produced by the region. This operation can be formulated without
discarding the depth history.

\begin{definition}[Admissibility of regional prolongation]
For fixed \(c,n\), let \(s=729^{n+1}\) and \(j=j_c(s)\). A residue
\(u\in\{0,\ldots,728\}\) is admissible when
\begin{equation}
 \lfloor s\ell_c(j)\rfloor
 \leq729P_c(729,n)+u
 \leq\lceil su_c(j)\rceil-1.
 \label{act21:eq:df-admisibilidad}
\end{equation}
\end{definition}

\begin{theorem}[Uniqueness of the prolongation compatible with the reader]
\label{act21:thm:df-prolongacion-unica}
Condition \eqref{act21:eq:df-admisibilidad} is equivalent to \(u=u_c(n)\).
There is therefore exactly one admissible prolongation of each region
at each depth.
\end{theorem}

\begin{proof}
At the separation index, both integer endpoints in
\eqref{act21:eq:df-admisibilidad} coincide with \(P_c(729,n+1)\).
The condition is therefore equivalent to
\(729P_c(729,n)+u=P_c(729,n+1)\).
By \eqref{act21:eq:df-paso-posicional}, the right-hand side is
\(729P_c(729,n)+u_c(n)\). Cancelling the common term gives equality
of the residues. Conversely, this equality satisfies the condition.
\end{proof}

The effective dependence of this selection is fully specified:
selected region, rational reader, depth, separation index, preceding
prefix, and next residue. The defect \(A\) summarizes a boundary
relationship; the state \((n,k,N,D)\) and the regional intervals retain
the data used by the update.

\subsection{Ternary signature, carry, Witt incidence, and reading clock}
\label{act21:subsec:df-firma-conjunta}

% Sources: lean/regional/N69RegionalSignature.lean,
% deltas/integracion/JointRegionalFrontier.lean and JointCylinderSignature.lean.

At a fixed depth \(t\), write \(B_{c,j}=B_{c,j}(t)\). For each column
\(j\), calculate the following integers and residues:
\begin{align}
 S_j&=B_{\mathrm P,j}+B_{\mathrm E,j}+B_{\mathrm F,j},
 &q_j&=S_j\bmod3,
 &c_j&=\left\lfloor S_j/3\right\rfloor,
 \label{act21:eq:df-firma-qc}\\
 a_j&=(B_{\mathrm P,j}+B_{\mathrm E,j}-B_{\mathrm F,j})\bmod3,
 &w_j&=\sum_c\mathbf1_{B_{c,j}\ne0},
 &\bar w_j&=w_j\bmod3.
 \label{act21:eq:df-firma-axial}
\end{align}
The residue of the axial expression is taken in the integers. An
implementation using natural numbers calculates it as
\((B_{\mathrm P,j}+B_{\mathrm E,j}+3-B_{\mathrm F,j})\bmod3\),
which preserves that residue exactly.

The ternary incidence matrix used by the reader is
\begin{equation}
 W=\begin{pmatrix}
0&1&1&1&1&1\\
1&0&1&1&2&2\\
1&1&0&2&1&2\\
1&1&2&0&2&1\\
1&2&1&2&0&1\\
1&2&2&1&1&0
\end{pmatrix}.
\label{act21:eq:df-matriz-witt}
\end{equation}
Two residual charges are preserved for each region:
\begin{equation}
 v_c^{\mathrm{res}}=\sum_j B_{c,j}\bmod3,
 \qquad
 v_c^{\mathrm W}=\sum_j(B_cW)_j\bmod3.
 \label{act21:eq:df-cargas-residuales}
\end{equation}
The reader's complete signature is
\(\mathcal S(B)=(q,a,c,w,\bar w,v^{\mathrm{res}},v^{\mathrm W})\).
Its structure separately preserves the column residue, carry, axial
contrast, support weight, and the two charge readings.

\begin{theorem}[Exact recovery and bounds of the signature]
\label{act21:thm:df-recuperacion-firma}
For each column and each region,
\begin{align}
 S_j&=q_j+3c_j,&0\leq q_j,a_j&<3,&0\leq c_j&\leq2,
 \label{act21:eq:df-recuperacion-total}\\
 B_{\mathrm F,j}&=2(q_j+3-a_j)\bmod3,&0\leq w_j&\leq3,
 &0\leq\bar w_j&<3.
 \label{act21:eq:df-recuperacion-eje}
\end{align}
Moreover,
\begin{equation}
 v_c^{\mathrm W}
 =\left(\sum_j\bigl((B_cW)_j\bmod3\bigr)\right)\bmod3.
 \label{act21:eq:df-witt-reduccion}
\end{equation}
The stated recoveries determine each column's total and its autoscaling
coordinate; the regional register additionally preserves the three
ordered bands.
\end{theorem}

\begin{proof}
Identity \eqref{act21:eq:df-recuperacion-total} is Euclidean division of
\(S_j\) by \(3\). Each summand of \(S_j\) lies between \(0\) and
\(2\), so \(0\leq S_j\leq6\) and \(c_j\leq2\). In \(\mathbb F_3\),
subtracting the two congruences defining \(q_j\) and \(a_j\) gives
\(q_j-a_j=2B_{\mathrm F,j}\). Since the inverse of \(2\) is \(2\),
this yields \eqref{act21:eq:df-recuperacion-eje}; the summand \(3\) permits
an expression using natural numbers. The bounds on \(w_j\) follow by
summing three indicators. Finally, reducing an integer sum modulo \(3\)
is equivalent to reducing each summand and then reducing their sum,
which proves \eqref{act21:eq:df-witt-reduccion}.
\end{proof}

\begin{example}[Signature of the first joint block]
The three initial bands produce
\begin{align*}
 q&=(0,0,2,2,1,2),&a&=(1,2,0,1,1,2),\\
 c&=(1,1,0,1,0,0),&w&=(2,2,2,3,1,2),\\
 v^{\mathrm{res}}&=(2,2,0),&v^{\mathrm W}&=(2,1,1).
\end{align*}
The fourth column has total \(2+1+2=5\); its signature records
\(q_3=2\) and \(c_3=1\), so \(q_3+3c_3=5\).
Axial recovery gives
\(2(2+3-1)\bmod3=2\), precisely the autoscaling coordinate.
The Witt images, reduced componentwise, are
\[
 (2,0,2,1,1,2),\qquad
 (0,0,0,2,0,2),\qquad
 (2,1,1,2,1,0).
\]
Their residual sums by row yield the three charges \((2,1,1)\).
\end{example}

\begin{definition}[Integer clock of decimal reading]
The transition index \(t\) and decimal expansion depth are related by
\begin{equation}
 \kappa(t)=\max\{k\in\mathbb N:1000^k\leq3^{6(t+1)}\}.
 \label{act21:eq:df-reloj}
\end{equation}
\end{definition}

\begin{proposition}[Characterization and monotonicity of the clock]
\label{act21:prop:df-reloj}
The integer \(\kappa(t)\) is uniquely determined by
\begin{equation}
 1000^{\kappa(t)}\leq729^{t+1}<1000^{\kappa(t)+1}.
 \label{act21:eq:df-reloj-cotas}
\end{equation}
The function \(\kappa\) is monotone and satisfies
\(\kappa(20)=\kappa(21)=20\).
\end{proposition}

\begin{proof}
The powers of \(1000\) are strictly increasing and unbounded; there is
therefore a unique interval between consecutive powers containing
\(729^{t+1}\). Monotonicity of the latter expression proves that of
\(\kappa\). The final two evaluations are the integer comparisons
\[
 1000^{20}\leq729^{21}<1000^{21},\qquad
 1000^{20}\leq729^{22}<1000^{21}.
\]
These comparisons use exact integer powers.
\end{proof}

Equality of decimal depth at two distinct transitions explains why both
indices must be retained: the reading clock can remain constant while
the ternary band advances and its signature is updated. Chronological
information resides in the ordered register of transitions.

\subsection{Composition theorem for the block, cylinder, and signature}

\begin{theorem}[Joint compatibility of regional readings]
\label{act21:thm:df-composicion-conjunta}
For each depth \(n\), there exists a unique triple of blocks
\((u_{\mathrm P},u_{\mathrm E},u_{\mathrm F})\) satisfying regional
admissibility \eqref{act21:eq:df-admisibilidad} in all three regions.
This triple is \((u_{\mathrm P}(n),u_{\mathrm E}(n),u_{\mathrm F}(n))\).
Simultaneously:
\begin{enumerate}
\item each block updates its cylinder through recurrences
\eqref{act21:eq:df-rec-ternaria} and \eqref{act21:eq:df-rec-conjunta};
\item its six trits produce the signature
\(\mathcal S(B_{\mathrm P}(n),B_{\mathrm E}(n),B_{\mathrm F}(n))\);
\item the column totals are recovered as \(q+3c\), and the Witt charges
are obtained from the same bands through
\eqref{act21:eq:df-witt-reduccion};
\item every deeper prefix returns exactly these blocks, truncated
cylinders, and signatures.
\end{enumerate}
\end{theorem}

\begin{proof}
Apply Theorem~\ref{act21:thm:df-prolongacion-unica} to each of the three
regional indices. Coordinatewise uniqueness gives joint uniqueness.
Substitute those same residues into the operations \(T_u,T_{u,d}\),
which allows Theorem~\ref{act21:thm:df-refinamiento-generado} to be applied.
The deterministic application of \eqref{act21:eq:df-bloque-y-banda} produces
the bands used by
\eqref{act21:eq:df-firma-qc}--\eqref{act21:eq:df-cargas-residuales}. The recovery
identities are those of Theorem~\ref{act21:thm:df-recuperacion-firma}.
Finally, Theorem~\ref{act21:thm:df-bloque-estable} and truncation in each base
prove stability at greater depth. At every step the identity of the
emitted block has been preserved; neither reading requires the selection
of a different block.
\end{proof}

The conclusion is an operational compatibility between cylindrical
contraction and preservation of incidence. Both use the same generated
register, retain the depths, and admit integer recovery. This result
is the explicit form of the connection
\[
 \text{emitted regional block}
 \longrightarrow
 \begin{cases}
 \text{cylinder and integer boundary},\\
 \text{band, signature, carry, and Witt charge}.
 \end{cases}
\]
The subsequent development of incidence and fields uses the second
reading while preserving its common provenance with the first.
